If @@M0@@, then @@M1@@ is surjective by 1.3 (ii); therefore @@M2@@ is surjective. Consequently, @@M3@@. On the other hand, if @@M4@@ is not surjective, then @@M5@@; therefore @@M6@@ is \uline{strictly} decreasing from @@M7@@ until it becomes @@M8@@; hence (iii).

\sgaStableTarget{sga6.e13.def.1-5}\textbf{Definition 1.5.} \emph{Let @@M9@@ be a coherent sheaf on @@M10@@, @@M11@@ an integer @@M12@@ and @@M13@@ a sequence of @@M14@@ integers. Then @@M15@@ is said to be a \emph{@@M16@@-sheaf} if it satisfies one of the following two equivalent sets of conditions}~:

\sgaStableTarget{sga6.e13.para.1-5-1}@@M17@@
\begin{enumerate}[label=(\roman*)]
\item \emph{At every point of the support of @@M18@@, @@M19@@ is generated by @@M20@@.}
\item \emph{@@M21@@.}
\item \emph{(When @@M22@@), there exists a section @@M23@@ defining an exact sequence @@M24@@ such that @@M25@@ is a @@M26@@-sheaf.}
\end{enumerate}

\sgaStableTarget{sga6.e13.para.1-5-2}@@M27@@ \emph{There exists an @@M28@@-regular sequence of @@M29@@ sections @@M30@@ such that, for @@M31@@, if @@M32@@ denotes the restriction of @@M33@@ to the zero scheme of @@M34@@, then @@M35@@.}

\sgaStableTarget{sga6.e13.prop.1-6}\textbf{Proposition 1.6.} \emph{Let @@M36@@ be a @@M37@@-sheaf on @@M38@@. Then}~:
\begin{enumerate}[label=(\roman*)]
\item \emph{For every ``sufficiently general'' sequence @@M39@@ of @@M40@@ sections @@M41@@, if @@M42@@ denotes the restriction of @@M43@@ to the zero scheme of @@M44@@, then @@M45@@, where @@M46@@.}
\item \emph{Every subsheaf @@M47@@ of @@M48@@ is a @@M49@@-sheaf.}
\end{enumerate}

